\subsection{Jets of maps}

12.1.1. Let \(x \in X\) and let \(f, g\) be two continuous maps defined in a neighborhood of \(x\) and with values in \(Y\). We say that \(f\) and \(g\) have contact of order \(\geqslant k\) at \(x\) if \(f(x) = g(x)\) and if there exist charts \((U, \varphi, E)\) of \(X\) at \(x\) and \((V, \psi, F)\) of \(Y\) at \(f(x)\) such that the maps \(\psi \circ f \circ \varphi^{-1}\) and \(\psi \circ g \circ \varphi^{-1}\), defined in a neighborhood of \(\varphi(x)\) in E and with values in F, have contact of order \(\geqslant k\) at \(\varphi(x)\) (1.1.2). This property is then verified by all charts of \(X\) at \(x\) and of \(Y\) at \(f(x)\).

12.1.2. Let \(x \in \mathbf{X}\), \(y \in \mathbf{Y}\). In the set of maps of class \(\mathbf{C}^r\) defined in a neighborhood of \(x\), with values in \(\mathbf{Y}\), and mapping \(x\) to \(y\), the relation « \(f\) and \(g\) have contact of order \(\geqslant k\) » is an equivalence relation; the class of \(f\) for this relation is denoted \(j_x^k(f)\) and is called the jet of order \(k\) of \(f\), with source \(x\) and target \(y\). The source of a jet \(j\) is denoted \(s(j)\) and its target \(b(j)\).

The set of jets of order k from X to Y (resp. with source x, resp. with target y) is denoted \(J^{k}(X, Y)\) (resp. \(J_{x}^{k}(X, Y)\) , resp. \(J^{k}(X, Y)_{y}\) ) and one sets

\[
\mathrm{J} _ {x} ^ {k} (\mathrm{X}, \mathrm{Y}) _ {y} = \mathrm{J} _ {x} ^ {k} (\mathrm{X}, \mathrm{Y}) \cap \mathrm{J} ^ {k} (\mathrm{X}, \mathrm{Y}) _ {y}.
\]

12.1.3. If U and V are open subsets of X and Y, respectively, we identify in the obvious way \(\mathrm{J}^{k}(\mathrm{U},\mathrm{V})\) with the inverse image of \(U \times V\) under the map

\[
(s, b): \mathrm{J} ^ {k} (\mathrm{X}, \mathrm{Y}) \rightarrow \mathrm{X} \times \mathrm{Y}.
\]

12.1.4. Let Z be a manifold of class \(C^{r}\), and let \((x, y, z) \in X \times Y \times Z\). Let \(j \in \mathrm{J}_{x}^{k}(X, Y)_{y}\) and let \(j' \in \mathrm{J}_{y}^{k}(Y, Z)_{z}\). The maps \(f' \circ f\), with \(f \in j\) and \(f' \in j'\), have the same kth-order jet at x; this jet is called the composite of j and \(j'\) and is denoted \(j' \circ j\); we have \(s(j' \circ j) = s(j)\) and \(b(j' \circ j) = b(j')\). If T is a manifold of class \(C^{r}\) and if \(j'' \in J_{z}^{k}(Z, T)\), one has

\[
j ^ {\prime \prime} \circ (j ^ {\prime} \circ j) = (j ^ {\prime \prime} \circ j ^ {\prime}) \circ j.
\]

12.1.5. Let \(j \in \mathrm{J}_x^k(\mathrm{X}, \mathrm{Y})\), with \(x \in \mathrm{X}\), and let \(k'\) be an integer such that \(0 \leqslant k' \leqslant k\). The jets \(j_x^{k'}(f)\), for \(f \in j\), are equal; the jet thus defined is denoted \(r^{k, k'}(j)\); the map \(r^{k, k'}: \mathrm{J}^k(\mathrm{X}, \mathrm{Y}) \to \mathrm{J}^{k'}(\mathrm{X}, \mathrm{Y})\) is surjective.

12.1.6. Suppose that \(r\) is equal to \(\infty\) or to \(\omega\). Let \(f\) and \(g\) be two maps of class \(C^{r}\) defined in a neighborhood of a point \(x \in X\) and with values in Y. If \(j_x^m(f) = j_x^m(g)\) for every integer \(m \geqslant 0\), one says that \(f\) and \(g\) have contact of infinite order at \(x\). One thus obtains an equivalence relation whose classes are called jets of infinite order at \(x\); the jet of infinite order of \(f\) is denoted \(j_x^\infty(f)\) or \(j_x^\omega(f)\). The definitions and results above extend without modification to jets of infinite order.

